\bookchapter{الفصل 29: الصحة الرقمية — إقامة علاقة صحية مع الهاتف}

\chapterintro{هل حدث لك: أردت تصفح الفيديوهات القصيرة لخمس دقائق، فرفعت رأسك بعد ساعتين؟ أو تصفح الهاتف ليلاً في السرير، يزداد نشاطاً، وتنام في الثانية فجراً؟ الهاتف غيّر حياتنا لكنه يبتلع وقتنا وصحتنا بهدوء. هذا الفصل لا يطلب منك “الإقلاع عن الهاتف”، بل يساعدك في استعادة السيطرة.}

\booksection{واحد: انظر “وقت الشاشة” أولاً — الأرقام توقظ}

قبل التغيير، افهم الوضع. في الهاتف وظيفة لم يفتحها كثيرون، تسمى “وقت استخدام الشاشة” (iPhone) أو “الصحة الرقمية” (Android)، تخبرك كم تنفق يومياً على الهاتف وأين.

\begin{itemize}
  \item \textbf{iPhone}: الإعدادات وقت استخدام الشاشة عرض كل النشاط. سترى عدد فتح القفل وعدد رفع الهاتف ومدة استخدام كل تطبيق.
  \item \textbf{أندرويد}: الإعدادات الصحة الرقمية والتحكم الأبوي. Huawei تسميه "استخدام صحي للهاتف"، Xiaomi "إدارة وقت الشاشة". وظائف مماثلة، جميعها تعرض البيانات اليومية.
\end{itemize}

\textbf{80\% ممن يرون هذه البيانات لأول مرة يُصدمون}. صديق اكتشف أنه ينفق أكثر من 4 ساعات يومياً على Douyin، 28 ساعة أسبوعياً — تعادل وظيفة بدوام جزئي. هذه “الصدمة الرقمية” الخطوة الأولى للتغيير.

\booksection{اثنان: ضع “حداً يومياً” للتطبيقات — ضع مقيّد سرعة لنفسك}

بعد معرفة المشكلة، اضبط قيوداً مستهدفة. لا تحذف التطبيق، فقط ضع “حصة يومية”:

\begin{enumerate}
  \item \textbf{iPhone}: الإعدادات وقت استخدام الشاشة حدود التطبيقات إضافة حد. اختر التطبيق (مثل Douyin)، اضبط 30 دقيقة يومياً. بعد الوقت تتحول أيقونة التطبيق للرمادي، ويحتاج إدخال كلمة مرور إضافية.
  \item \textbf{أندرويد}: الإعدادات الصحة الرقمية مؤقت التطبيقات. اختر التطبيق واضبط الحد اليومي، بعد الوقت تتحول الأيقونة للرمادي.
  \item \textbf{نصيحة صغيرة}: اجعل كلمة مرور الحد كلمة غير شائعة وصعبة الحفظ. إذا كانت سهلة، ستضيف وقتاً في كل مرة، ويصبح الحد بلا معنى.
\end{enumerate}

\booksection{ثلاثة: الاستخدام الصحيح لـ “وضعدم الإزعاج” — لا إغلاق الإشعارات بل إغلاق القلق}

هاتفك يرن عشرات المرات يومياً، كل رنة تقطع انتباهك. تظن أنك “تنظر إشعاراً”، لكن إعادة التركيز بعد كل مقاطعة تحتاج 15-20 دقيقة. وضععدم الإزعاج لا يجعلك “منفصلاً” بل يعيد السيطرة على انتباهك.

\begin{itemize}
  \item \textbf{تفعيل مجدول}: اضبط وضععدم الإزعاج للتفعيل التلقائي من 10 مساءً إلى 7 صباحاً. هكذا لا يوقظك هاتفك برسالة مزعجة.
  \item \textbf{السماح بمرور جهات الاتصال المهمة}: في إعدادات وضععدم الإزعاج، اضبط "المكالمات المسموحة" لـ "جهات الاتصال المفضلة" أو "المكالمات المتكررة" (المكالمة الثانية خلال 3 دقائق ترن). هكذا لا تفوت مكالمة طوارئ من الأسرة، دون إزعاج من مكالمات البيع.
  \item \textbf{وضع تركيز العمل}: "وضع التركيز" في iPhone يمكن ضبط سيناريوهات مخصصة — مثلاً وضع "العمل" يسمح فقط بإشعارات التطبيقات المتعلقة بالعمل، وضع "القراءة" يكتم جميع الإشعارات.
\end{itemize}

\booksection{أربعة: حماية العين بالعادة لا بالتطبيق — قاعدة 20-20-20}

النظر الطويل للشاشة يجفف العين ويتعبها. وضع حماية العين (فلتر الضوء الأزرق) يساعد بعض الشيء، لكن الأهم هو تغيير عادات العين:

\begin{tcbinfo}{قاعدة 20-20-20}
كل\textbf{20 دقيقة}نظر للشاشة، انقل نظرك لـ\textbf{20 قدماً (حوالي 6 أمتار)}بعيداً، لمدة\textbf{20 ثانية}. هذا الإجراء البسيط يخفف إجهاد العضلة الهدبية بفعالية.

\end{tcbinfo}

\begin{itemize}
  \item \textbf{تفعيل وضع حماية العين}: الإعدادات العرض وضع حماية العين (أو الوضع الليلي)، بعد التفعيل يصبح لون الشاشة دافئاً، يقلل الضوء الأزرق. يُنصح بضبطه للتفعيل التلقائي من الغروب للشروق.
  \item \textbf{ضبط سطوع الشاشة}: في البيئة المظلمة يجب أن تكون الشاشة مظلمة. كثيرون يضبطون السطوع عالياً في الظلام، وهذا أضرر سلوك للعين. فعّل "السطوع التلقائي" وتأكد من عمله.
  \item \textbf{الحفاظ على المسافة}: الهاتف على بُعد 30 سم على الأقل من العين. كثيرون يقرّبون الهاتف جداً، هكذا تبقى العضلة الهدبية متوترة، ويزداد قصر النظر بسرعة.
\end{itemize}

\booksection{خمسة: دليل إنقاذ الرقبة “للمنخفضين”}

عند الانحناء للهاتف، يتحمل رقبتك ضغطاً 5-6 أضعاف الوقوف العادي. الانحناء الطويل يؤدي لتسطيح تقوس الرقبة وتنكس العظام ودوخة وتنميل اليد. التالي طرق تحسن بسيطة:

\begin{itemize}
  \item \textbf{ارفع الهاتف}: عند استخدام الهاتف ارفعه لمستوى العين، بدلاً من الانحناء للهاتف على فخذيك. ذراعك ستتعب، لكن رقبتك ستشكرك.
  \item \textbf{حرّك رقبتك كل 30 دقيقة}: أدر رأسك ببطء لليسار للنهاية، توقف 3 ثوانٍ، عُد؛ ثم لليمين، توقف 3 ثوانٍ. كرر 5 مرات. ثم اعمل "تمرين حرف مي" — اكتب بحركة رقبتك في الهاء حرف "Mi".
  \item \textbf{الوقوف على الحائط 3 دقائق}: مؤخرة الرأس واللوحان الكتفيان والوركان والكعبان أربع نقاط تلامس الحائط، قف 3 دقائق يومياً، يصحح وضعية الانحناء بفعالية.
\end{itemize}

\booksection{ستة: النوم والهاتف — العادة التي يجب تغييرها في غرفة النوم}

الضوء الأزرق من شاشة الهاتف يثبط إفراز الميلاتونين، يجعلك “تزداد نشاطاً”. كثيرون يظنون “أتصفح فينعسني”، لكن في الواقع الضوء الأزرق يخدع الدماغ، وعند وضع الهاتف يصعب النوم. اقتراحات:

\begin{itemize}
  \item \textbf{لا تنظر للهاتف قبل النوم بساعة}: هذا الأكثر فعالية والأصعب. إن لم تستطع، فعّل على الأقل "الوضع الليلي/حماية العين" لأدنى ضوء أزرق.
  \item \textbf{ضع الهاتف بعيداً عن متناول اليد}: ضع الشاحن في الطرف الآخر من الغرفة، ليس على طاولة السرير. هكذا لن تمسكه، وعند رنين المنبه عليك النهوض لإيقافه — فائدتان.
  \item \textbf{تفعيل "وضع النوم"}: "وضع تركيز النوم" في iPhone يقفل الشاشة تلقائياً في وقت النوم ويبسّط شاشة القفل؛ Android لكل علامة وظائف مماثلة.
  \item \textbf{استخدم منبهاً بدلاً من الهاتف}: اشترِ منبهاً مادياً بعشرة يوانات، لا تستخدم الهاتف كمنبه، ستكتشف اختفاء فخ "أنظر للوقت ثم أتصفح".
\end{itemize}

\booksection{سبعة: قلق المعلومات — تعلم “العمى الاختياري”}

كمية المعلومات التي نتلقاها يومياً تعادل ما يتلقاه قديم طوال حياته. لكن دماغ الإنسان لم يُصمم لمعالجة هذا الكم، التحميل الزائد يسبب القلق وتشتت الانتباه وإرهاق القرار. الحل ليس “الحصول على المزيد” بل “الحجب النشط”:

\begin{itemize}
  \item \textbf{إلغاء المتابعة}: افتح قائمة الحسابات العامة وUP والمؤثرين الذين تتابعهم، اسأل نفسك: "هل حصلت على قيمة فعلية في الشهر الماضي؟" إن لا، ألغِ المتابعة. القنوات النوعية لا الكمية.
  \item \textbf{إغلاق إشعارات التطبيقات غير الضرورية}: باستثناء WeChat والهاتف والرسائل والتقويم، أغلق إشعارات جميع التطبيقات. الأخبار لا تحتاج إشعارات — ستراها عند البحث بنفسك.
  \item \textbf{تأسيس "يوم صيام المعلومات"}: اختر يوماً أسبوعياً (مثل الأحد)، لا تنظر لأخبار الهاتف ولا تتصفح الفيديوهات القصيرة. ستكتشف أن السماء لا تسقط، وحالتك يوم الاثنين أفضل بوضوح.
\end{itemize}

\begin{tcbpractice}{جرّب: تحدي الصحة الرقمية لأسبوع}
\begin{enumerate}
  \item اليوم: افتح وقت استخدام الشاشة/الصحة الرقمية، راجع بيانات الأسبوع الماضي.
  \item غداً: اضبط حداً يومياً للتطبيقات الترفيهية الأكثر استهلاكاً للوقت (يُنصح بـ 30-60 دقيقة).
  \item بعد غد: اضبط وضععدم الإزعاج للتفعيل المجدول (22:00-07:00)، اسمح بمكالمات جهات الاتصال المفضلة.
  \item اليوم الرابع: لا تنظر للهاتف قبل النوم بساعة، ضع الشاحن في الطرف الآخر من الغرفة.
  \item اليوم الخامس: ألغِ متابعة 5 حسابات أو مؤثرين لم تعد تتابعهم.
  \item اليوم السادس: جرّب "صيام المعلومات" ليوم كامل — لا أخبار ولا فيديوهات قصيرة.
  \item اليوم السابع: راجع تقرير وقت استخدام الشاشة لهذا الأسبوع، وقارن بالأسبوع الماضي.
\end{enumerate}

\end{tcbpractice}

\begin{tcbsummary}{خلاصة الفصل}
\begin{itemize}
  \item راجع بيانات وقت الشاشة أولاً، إدراك الواقع هو الخطوة الأولى للتغيير
  \item ضع حدوداً للتطبيقات الترفيهية، واضبط وضععدم الإزعاج، استعد السيطرة على الانتباه
  \item قاعدة 20-20-20 أبسط وأ فعالية عادة لحماية العين
  \item الانحناء للهاتف يحمّل الرقبة 5-6 أضعاف، ارفع الهاتف + حرّك رقبتك دورياً
  \item لا تنظر للهاتف قبل النوم بساعة، ضع الشاحن بعيداً عن متناول اليد
  \item دواء قلق المعلومات "الحجب النشط"، لا "الحصول على المزيد"
\end{itemize}

\end{tcbsummary}
